% --- Archon physics formalization source begin ---
% archon:physics
% archon:covers PhyXMiniProblems/problem_phyx_mini_0599.lean
% archon:source-report reports/phyx_mini/problem_phyx_mini_0599.source.json

\chapter{Physics problem phyx\_mini\_0599}
\label{ch:PhyXMiniProblems_problem_phyx_mini_0599}

\paragraph{Problem source.}
\#\# Physical scenario

Figure gives \textbackslash{}( \textbackslash{}Delta t \textbackslash{}) versus speed parameter \textbackslash{}( eta \textbackslash{}) for a range of values for \textbackslash{}( eta \textbackslash{}). The vertical axis scale is set by \textbackslash{}( \textbackslash{}Delta t\_a = 14.0\textbackslash{},\textbackslash{}text\{s\} \textbackslash{}).

\#\# Question

What is interval \textbackslash{}( \textbackslash{}Delta t \textbackslash{}) if \textbackslash{}( v = 0.98c \textbackslash{})?

\#\# Answer choices

- A: A: 28s

- B: B: 56s

- C: C: 40s

- D: D: 12s

\#\# Figure caption (auxiliary; use the image as primary evidence)

The image shows a graph with a thick curve plotting the relationship between two variables, \textbackslash{}(\textbackslash{}beta\textbackslash{}) and \textbackslash{}(\textbackslash{}Delta t\_a\textbackslash{}). 

- **X-axis**: Labeled as \textbackslash{}(\textbackslash{}beta\textbackslash{}), with values ranging from 0 to 0.8, marked at intervals of 0.4.
- **Y-axis**: Labeled as \textbackslash{}(\textbackslash{}Delta t (s)\textbackslash{}) with \textbackslash{}(\textbackslash{}Delta t\_a\textbackslash{}) noted at the top, divided into equal sections without specific values.
- **Line/Curve**: The curve starts at the origin (0,0) and increases non-linearly, becoming steeper as \textbackslash{}(\textbackslash{}beta\textbackslash{}) increases.

The graph is overlaid with a grid for easier reading, and the font used is likely a serif style, typical for mathematical and scientific texts.

\#\# Dataset metadata

Relativity; Physical Model Grounding Reasoning, Multi-Formula Reasoning

\paragraph{Recorded answer/context.}
C: C: 40s

\paragraph{Figure/image path.}
/root/proposal\_for\_physic/science-mango/phyx\_mini\_run/phyx\_data/test\_image/599.png

\paragraph{Formalization target.}
create a compiling Lean file with sorry bodies at `PhyXMiniProblems/problem\_phyx\_mini\_0599.lean`. The Lean declarations must preserve the physical quantities, dimensions or dimensional roles, figure labels, governing-law hypotheses, and final relation expressed by this problem.
Use Mathlib/Physlib names found through LeanExplore where available. If a physics API is missing, introduce faithful local abstractions rather than scalar placeholder aliases.

\begin{theorem}[Physics formalization target]
\label{thm:physics:phyx_mini_0599:target}
\lean{PhyXMiniProblems.ProblemPhyXMini0599.problem_phyx_mini_0599}
\uses{def:physics:phyx-mini-0599:phyxminiproblems-problemphyxmini0599-durationinseconds, def:physics:phyx-mini-0599:phyxminiproblems-problemphyxmini0599-relativisticintervalsetup, def:physics:phyx-mini-0599:phyxminiproblems-problemphyxmini0599-speedfractionoflight, def:physics:phyx-mini-0599:phyxminiproblems-problemphyxmini0599-matchestimedilationscenario, def:physics:phyx-mini-0599:phyxminiproblems-problemphyxmini0599-matchessuppliedtimedilationfigure, def:physics:phyx-mini-0599:phyxminiproblems-problemphyxmini0599-matchesproblemreadouts, def:physics:phyx-mini-0599:phyxminiproblems-problemphyxmini0599-hasphysicalrelativisticparameters, def:physics:phyx-mini-0599:phyxminiproblems-problemphyxmini0599-satisfieslorentztimedilation, def:physics:phyx-mini-0599:phyxminiproblems-problemphyxmini0599-recordedanswerchoice, def:physics:phyx-mini-0599:phyxminiproblems-problemphyxmini0599-matchesdisplayedtimechoice, def:physics:phyx-mini-0599:phyxminiproblems-problemphyxmini0599-isuniqueclosesttimechoice}
The assigned autoformalize agent should translate this physics problem into Lean declarations in the covered file, with theorem and lemma proof bodies written as `by sorry`.
\end{theorem}
\begin{proof}
This is an autoformalization task, not a proof task. Produce statements that can later be proved by the physics prover without weakening the physical model.
\end{proof}

% --- Archon declaration topology begin ---
\section{Declaration topology}
\begin{definition}[Duration Quantity]
  \label{def:physics:phyx-mini-0599:phyxminiproblems-problemphyxmini0599-durationquantity}
  \lean{PhyXMiniProblems.ProblemPhyXMini0599.DurationQuantity}
  A nonnegative, unit-independent physical duration
\end{definition}
\begin{proof}
  This is the declared model definition of Duration Quantity.
\end{proof}
\begin{definition}[Speed Quantity]
  \label{def:physics:phyx-mini-0599:phyxminiproblems-problemphyxmini0599-speedquantity}
  \lean{PhyXMiniProblems.ProblemPhyXMini0599.SpeedQuantity}
  A physical speed carrying length-per-time dimension
\end{definition}
\begin{proof}
  This is the declared model definition of Speed Quantity.
\end{proof}
\begin{definition}[duration Readout]
  \label{def:physics:phyx-mini-0599:phyxminiproblems-problemphyxmini0599-durationreadout}
  \lean{PhyXMiniProblems.ProblemPhyXMini0599.durationReadout}
  \uses{def:physics:phyx-mini-0599:phyxminiproblems-problemphyxmini0599-durationquantity}
  Read a physical duration in a selected time unit
\end{definition}
\begin{proof}
  This is the declared model definition of duration Readout.
\end{proof}
\begin{definition}[duration In Seconds]
  \label{def:physics:phyx-mini-0599:phyxminiproblems-problemphyxmini0599-durationinseconds}
  \lean{PhyXMiniProblems.ProblemPhyXMini0599.durationInSeconds}
  \uses{def:physics:phyx-mini-0599:phyxminiproblems-problemphyxmini0599-durationquantity, def:physics:phyx-mini-0599:phyxminiproblems-problemphyxmini0599-durationreadout}
  Read a physical duration in seconds, as on the vertical axis
\end{definition}
\begin{proof}
  This is the declared model definition of duration In Seconds.
\end{proof}
\begin{definition}[speed Readout]
  \label{def:physics:phyx-mini-0599:phyxminiproblems-problemphyxmini0599-speedreadout}
  \lean{PhyXMiniProblems.ProblemPhyXMini0599.speedReadout}
  \uses{def:physics:phyx-mini-0599:phyxminiproblems-problemphyxmini0599-speedquantity}
  Read a physical speed in coherent selected length and time units
\end{definition}
\begin{proof}
  This is the declared model definition of speed Readout.
\end{proof}
\begin{definition}[speed In Meters Per Second]
  \label{def:physics:phyx-mini-0599:phyxminiproblems-problemphyxmini0599-speedinmeterspersecond}
  \lean{PhyXMiniProblems.ProblemPhyXMini0599.speedInMetersPerSecond}
  \uses{def:physics:phyx-mini-0599:phyxminiproblems-problemphyxmini0599-speedquantity, def:physics:phyx-mini-0599:phyxminiproblems-problemphyxmini0599-speedreadout}
  Read a physical speed in metres per second
\end{definition}
\begin{proof}
  This is the declared model definition of speed In Meters Per Second.
\end{proof}
\begin{definition}[vacuum Speed Of Light In Meters Per Second]
  \label{def:physics:phyx-mini-0599:phyxminiproblems-problemphyxmini0599-vacuumspeedoflightinmeterspersecond}
  \lean{PhyXMiniProblems.ProblemPhyXMini0599.vacuumSpeedOfLightInMetersPerSecond}
  Physlib's exact vacuum speed of light in metres per second
\end{definition}
\begin{proof}
  This is the declared model definition of vacuum Speed Of Light In Meters Per Second.
\end{proof}
\begin{definition}[Inertial Frame Label]
  \label{def:physics:phyx-mini-0599:phyxminiproblems-problemphyxmini0599-inertialframelabel}
  \lean{PhyXMiniProblems.ProblemPhyXMini0599.InertialFrameLabel}
  The rest frame of the clock and the frame in which Delta t is plotted
\end{definition}
\begin{proof}
  This is the declared model definition of Inertial Frame Label.
\end{proof}
\begin{definition}[Figure Axis]
  \label{def:physics:phyx-mini-0599:phyxminiproblems-problemphyxmini0599-figureaxis}
  \lean{PhyXMiniProblems.ProblemPhyXMini0599.FigureAxis}
  The two axes visible in the supplied graph
\end{definition}
\begin{proof}
  This is the declared model definition of Figure Axis.
\end{proof}
\begin{definition}[Figure Axis Label]
  \label{def:physics:phyx-mini-0599:phyxminiproblems-problemphyxmini0599-figureaxislabel}
  \lean{PhyXMiniProblems.ProblemPhyXMini0599.FigureAxisLabel}
  Literal mathematical labels printed on the graph
\end{definition}
\begin{proof}
  This is the declared model definition of Figure Axis Label.
\end{proof}
\begin{definition}[Figure Scale Label]
  \label{def:physics:phyx-mini-0599:phyxminiproblems-problemphyxmini0599-figurescalelabel}
  \lean{PhyXMiniProblems.ProblemPhyXMini0599.FigureScaleLabel}
  The scale annotation printed at the upper end of the vertical axis
\end{definition}
\begin{proof}
  This is the declared model definition of Figure Scale Label.
\end{proof}
\begin{definition}[Time Dilation Figure]
  \label{def:physics:phyx-mini-0599:phyxminiproblems-problemphyxmini0599-timedilationfigure}
  \lean{PhyXMiniProblems.ProblemPhyXMini0599.TimeDilationFigure}
  \uses{def:physics:phyx-mini-0599:phyxminiproblems-problemphyxmini0599-figureaxis, def:physics:phyx-mini-0599:phyxminiproblems-problemphyxmini0599-figureaxislabel, def:physics:phyx-mini-0599:phyxminiproblems-problemphyxmini0599-figurescalelabel}
  Literal information carried by phyx\_data/test\_image/599.png. normalizedCurveHeight is a dimensionless drawing coordinate: 0 is the bottom axis and 1 is the horizontal line carrying the Delta t\_a scale. It is not itself a physical duration
\end{definition}
\begin{proof}
  This is the declared model definition of Time Dilation Figure.
\end{proof}
\begin{definition}[Relativistic Interval Setup]
  \label{def:physics:phyx-mini-0599:phyxminiproblems-problemphyxmini0599-relativisticintervalsetup}
  \lean{PhyXMiniProblems.ProblemPhyXMini0599.RelativisticIntervalSetup}
  \uses{def:physics:phyx-mini-0599:phyxminiproblems-problemphyxmini0599-durationquantity, def:physics:phyx-mini-0599:phyxminiproblems-problemphyxmini0599-speedquantity, def:physics:phyx-mini-0599:phyxminiproblems-problemphyxmini0599-inertialframelabel, def:physics:phyx-mini-0599:phyxminiproblems-problemphyxmini0599-timedilationfigure}
  The coordinate interval at each speed parameter, proper interval, axis-scale interval, and stated physical speed are independent quantities. In particular, no field is defined from the recorded 40 s answer
\end{definition}
\begin{proof}
  This is the declared model definition of Relativistic Interval Setup.
\end{proof}
\begin{definition}[speed Fraction Of Light]
  \label{def:physics:phyx-mini-0599:phyxminiproblems-problemphyxmini0599-speedfractionoflight}
  \lean{PhyXMiniProblems.ProblemPhyXMini0599.speedFractionOfLight}
  \uses{def:physics:phyx-mini-0599:phyxminiproblems-problemphyxmini0599-speedinmeterspersecond, def:physics:phyx-mini-0599:phyxminiproblems-problemphyxmini0599-vacuumspeedoflightinmeterspersecond, def:physics:phyx-mini-0599:phyxminiproblems-problemphyxmini0599-relativisticintervalsetup}
  The dimensionless speed parameter beta = v/c in coherent SI readouts
\end{definition}
\begin{proof}
  This is the declared model definition of speed Fraction Of Light.
\end{proof}
\begin{definition}[Matches Time Dilation Scenario]
  \label{def:physics:phyx-mini-0599:phyxminiproblems-problemphyxmini0599-matchestimedilationscenario}
  \lean{PhyXMiniProblems.ProblemPhyXMini0599.MatchesTimeDilationScenario}
  \uses{def:physics:phyx-mini-0599:phyxminiproblems-problemphyxmini0599-relativisticintervalsetup}
  The two duration roles belong to the inertial frames used by time dilation
\end{definition}
\begin{proof}
  This is the declared model definition of Matches Time Dilation Scenario.
\end{proof}
\begin{definition}[Matches Supplied Time Dilation Figure]
  \label{def:physics:phyx-mini-0599:phyxminiproblems-problemphyxmini0599-matchessuppliedtimedilationfigure}
  \lean{PhyXMiniProblems.ProblemPhyXMini0599.MatchesSuppliedTimeDilationFigure}
  \uses{def:physics:phyx-mini-0599:phyxminiproblems-problemphyxmini0599-timedilationfigure}
  Primary-image evidence. The raster, rather than the auxiliary caption, shows the curve at four sevenths of the vertical scale when beta = 0
\end{definition}
\begin{proof}
  This is the declared model definition of Matches Supplied Time Dilation Figure.
\end{proof}
\begin{definition}[Matches Problem Readouts]
  \label{def:physics:phyx-mini-0599:phyxminiproblems-problemphyxmini0599-matchesproblemreadouts}
  \lean{PhyXMiniProblems.ProblemPhyXMini0599.MatchesProblemReadouts}
  \uses{def:physics:phyx-mini-0599:phyxminiproblems-problemphyxmini0599-durationinseconds, def:physics:phyx-mini-0599:phyxminiproblems-problemphyxmini0599-relativisticintervalsetup, def:physics:phyx-mini-0599:phyxminiproblems-problemphyxmini0599-speedfractionoflight}
  The 14 s scale and v = 0.98 c are stated data. On the displayed beta range, the curve height is connected to the corresponding dimensionful coordinate interval. This relation contains no readout at the queried beta = 0.98, which lies outside the graph's displayed range
\end{definition}
\begin{proof}
  This is the declared model definition of Matches Problem Readouts.
\end{proof}
\begin{definition}[Has Physical Relativistic Parameters]
  \label{def:physics:phyx-mini-0599:phyxminiproblems-problemphyxmini0599-hasphysicalrelativisticparameters}
  \lean{PhyXMiniProblems.ProblemPhyXMini0599.HasPhysicalRelativisticParameters}
  \uses{def:physics:phyx-mini-0599:phyxminiproblems-problemphyxmini0599-durationinseconds, def:physics:phyx-mini-0599:phyxminiproblems-problemphyxmini0599-vacuumspeedoflightinmeterspersecond, def:physics:phyx-mini-0599:phyxminiproblems-problemphyxmini0599-relativisticintervalsetup, def:physics:phyx-mini-0599:phyxminiproblems-problemphyxmini0599-speedfractionoflight}
  Positivity and the nonnegative subluminal regime required by the model
\end{definition}
\begin{proof}
  This is the declared model definition of Has Physical Relativistic Parameters.
\end{proof}
\begin{definition}[Satisfies Lorentz Time Dilation]
  \label{def:physics:phyx-mini-0599:phyxminiproblems-problemphyxmini0599-satisfieslorentztimedilation}
  \lean{PhyXMiniProblems.ProblemPhyXMini0599.SatisfiesLorentzTimeDilation}
  \uses{def:physics:phyx-mini-0599:phyxminiproblems-problemphyxmini0599-durationinseconds, def:physics:phyx-mini-0599:phyxminiproblems-problemphyxmini0599-relativisticintervalsetup}
  The governing special-relativistic law. A coordinate interval at speed parameter beta equals the proper interval multiplied by Physlib's Lorentz factor. It is stated for arbitrary nonnegative subluminal beta, rather than only for the numerical target of this problem
\end{definition}
\begin{proof}
  This is the declared model definition of Satisfies Lorentz Time Dilation.
\end{proof}
\begin{lemma}[proper Interval seconds eq eight]
  \label{lem:physics:phyx-mini-0599:phyxminiproblems-problemphyxmini0599-properinterval-seconds-eq-eight}
  \lean{PhyXMiniProblems.ProblemPhyXMini0599.properInterval_seconds_eq_eight}
  \uses{def:physics:phyx-mini-0599:phyxminiproblems-problemphyxmini0599-durationinseconds, def:physics:phyx-mini-0599:phyxminiproblems-problemphyxmini0599-relativisticintervalsetup, def:physics:phyx-mini-0599:phyxminiproblems-problemphyxmini0599-matchessuppliedtimedilationfigure, def:physics:phyx-mini-0599:phyxminiproblems-problemphyxmini0599-matchesproblemreadouts, def:physics:phyx-mini-0599:phyxminiproblems-problemphyxmini0599-hasphysicalrelativisticparameters, def:physics:phyx-mini-0599:phyxminiproblems-problemphyxmini0599-satisfieslorentztimedilation}
  Four of the seven equal vertical divisions of a 14 s scale calibrate the curve at beta = 0 to 8 s. Since gamma(0) = 1, this is also the proper interval. The value is a consequence of the image and law, not a premise
\end{lemma}
\begin{proof}
  The conclusion follows from the governing relations and figure readouts named in the statement of proper Interval seconds eq eight.
\end{proof}
\begin{lemma}[interval At Stated Speed exact Formula]
  \label{lem:physics:phyx-mini-0599:phyxminiproblems-problemphyxmini0599-intervalatstatedspeed-exactformula}
  \lean{PhyXMiniProblems.ProblemPhyXMini0599.intervalAtStatedSpeed_exactFormula}
  \uses{def:physics:phyx-mini-0599:phyxminiproblems-problemphyxmini0599-durationinseconds, def:physics:phyx-mini-0599:phyxminiproblems-problemphyxmini0599-relativisticintervalsetup, def:physics:phyx-mini-0599:phyxminiproblems-problemphyxmini0599-speedfractionoflight, def:physics:phyx-mini-0599:phyxminiproblems-problemphyxmini0599-matchessuppliedtimedilationfigure, def:physics:phyx-mini-0599:phyxminiproblems-problemphyxmini0599-matchesproblemreadouts, def:physics:phyx-mini-0599:phyxminiproblems-problemphyxmini0599-hasphysicalrelativisticparameters, def:physics:phyx-mini-0599:phyxminiproblems-problemphyxmini0599-satisfieslorentztimedilation}
  At the stated speed, the observer-frame interval is the calibrated 8 s proper interval multiplied by Physlib's Lorentz factor at 49/50
\end{lemma}
\begin{proof}
  The conclusion follows from the governing relations and figure readouts named in the statement of interval At Stated Speed exact Formula.
\end{proof}
\begin{definition}[Answer Choice]
  \label{def:physics:phyx-mini-0599:phyxminiproblems-problemphyxmini0599-answerchoice}
  \lean{PhyXMiniProblems.ProblemPhyXMini0599.AnswerChoice}
  Labels of the four whole-second choices printed with the problem
\end{definition}
\begin{proof}
  This is the declared model definition of Answer Choice.
\end{proof}
\begin{definition}[seconds]
  \label{def:physics:phyx-mini-0599:phyxminiproblems-problemphyxmini0599-answerchoice-seconds}
  \lean{PhyXMiniProblems.ProblemPhyXMini0599.AnswerChoice.seconds}
  \uses{def:physics:phyx-mini-0599:phyxminiproblems-problemphyxmini0599-answerchoice}
  Seconds printed beside each answer label
\end{definition}
\begin{proof}
  This is the declared model definition of seconds.
\end{proof}
\begin{definition}[recorded Answer Choice]
  \label{def:physics:phyx-mini-0599:phyxminiproblems-problemphyxmini0599-recordedanswerchoice}
  \lean{PhyXMiniProblems.ProblemPhyXMini0599.recordedAnswerChoice}
  \uses{def:physics:phyx-mini-0599:phyxminiproblems-problemphyxmini0599-answerchoice}
  The answer label recorded by the source dataset
\end{definition}
\begin{proof}
  This is the declared model definition of recorded Answer Choice.
\end{proof}
\begin{definition}[Matches Displayed Time Choice]
  \label{def:physics:phyx-mini-0599:phyxminiproblems-problemphyxmini0599-matchesdisplayedtimechoice}
  \lean{PhyXMiniProblems.ProblemPhyXMini0599.MatchesDisplayedTimeChoice}
  \uses{def:physics:phyx-mini-0599:phyxminiproblems-problemphyxmini0599-durationquantity, def:physics:phyx-mini-0599:phyxminiproblems-problemphyxmini0599-durationinseconds, def:physics:phyx-mini-0599:phyxminiproblems-problemphyxmini0599-answerchoice}
  Agreement with a displayed whole-second value after nearest-second rounding
\end{definition}
\begin{proof}
  This is the declared model definition of Matches Displayed Time Choice.
\end{proof}
\begin{definition}[Is Unique Closest Time Choice]
  \label{def:physics:phyx-mini-0599:phyxminiproblems-problemphyxmini0599-isuniqueclosesttimechoice}
  \lean{PhyXMiniProblems.ProblemPhyXMini0599.IsUniqueClosestTimeChoice}
  \uses{def:physics:phyx-mini-0599:phyxminiproblems-problemphyxmini0599-durationquantity, def:physics:phyx-mini-0599:phyxminiproblems-problemphyxmini0599-durationinseconds, def:physics:phyx-mini-0599:phyxminiproblems-problemphyxmini0599-answerchoice}
  The chosen value is strictly closer than every alternative displayed value
\end{definition}
\begin{proof}
  This is the declared model definition of Is Unique Closest Time Choice.
\end{proof}
% --- Archon declaration topology end ---
% --- Archon physics formalization source end ---
